% unit: u08
% sources: sources/014-solving-ax-0-pivot-variables-special-solutions.txt, sources/024-solving-ax-0-pivot-variables-special-solutions.txt, sources/107-mit18-06scf11-ses1-7sum-pdf-download-file.txt, sources/108-solving-ax-0.txt, sources/109-mit18-06scf11-ses1-7prob-pdf-download-file.txt, sources/110-mit18-06scf11-ses1-7sol-pdf-download-file.txt
\section{\texorpdfstring{Solving \(A\vb{x}=\vb{0}\): pivot variables, special solutions}{Solving Ax = 0: pivot variables, special solutions}}\label{sec:u08}

Section~\ref{sec:u07} defined the nullspace $\Nul{A}$ as the set of all solutions of
$A\vb{x}=\vb{0}$ and showed that it is a subspace of $\R^n$. A definition is not yet a
description: we want to \emph{list} what is in $\Nul{A}$. This section turns elimination, which
so far solved square invertible systems, into an algorithm that works for every $m\times n$
matrix. The algorithm produces one number, the rank $r$, and a short list of $n-r$
\emph{special solutions} whose combinations are exactly the nullspace.

\subsection{Elimination on a rectangular matrix}

The running example of the lecture is the $3\times 4$ matrix
\[
A=\begin{bmatrix}1&2&2&2\\2&4&6&8\\3&6&8&10\end{bmatrix}.
\]
Two facts are visible before any computation: column~2 is twice column~1, and row~3 is
row~1 plus row~2. Elimination should discover both. The first question is whether elimination
is allowed to change the problem.

\begin{proposition}[Row operations keep the nullspace]\label{prop:u08-same-null}
If $E$ is an invertible $m\times m$ matrix, then $\Nul{EA}=\Nul{A}$. In particular, subtracting
a multiple of one row from another, exchanging two rows, or dividing a row by a nonzero number
never changes the nullspace.
\end{proposition}
\begin{proof}
If $A\vb{x}=\vb{0}$ then $EA\vb{x}=E\vb{0}=\vb{0}$. Conversely, if $EA\vb{x}=\vb{0}$, multiply on
the left by $E^{-1}$ to get $A\vb{x}=\vb{0}$. Each elimination step is multiplication on the left
by an invertible elementary or permutation matrix (Sections~\ref{sec:u03} and~\ref{sec:u06}), and
a product of invertible matrices is invertible.
\end{proof}

Because the right-hand side is $\vb{0}$ and stays $\vb{0}$ under every row operation, we do not
need to carry it along. Elimination proceeds column by column as in the square case, with one
new rule: if the pivot position and every entry below it are zero, there is nothing to do in
that column, so we move to the next column (staying in the same row) instead of stopping.

\begin{example}[Elimination with a missing pivot]\label{ex:u08-elim}
Subtract $2\times$row~1 from row~2 and $3\times$row~1 from row~3:
\[
\begin{bmatrix}1&2&2&2\\2&4&6&8\\3&6&8&10\end{bmatrix}\longrightarrow
\begin{bmatrix}1&2&2&2\\0&0&2&4\\0&0&2&4\end{bmatrix}.
\]
In column~2 the pivot position $(2,2)$ holds a zero and so does the entry below it, so no row
exchange can help: column~2 gets no pivot. The next pivot is the $2$ in row~2, column~3.
Subtracting row~2 from row~3 finishes the job:
\[
U=\begin{bmatrix}\boxed{1}&2&2&2\\0&0&\boxed{2}&4\\0&0&0&0\end{bmatrix}.
\]
The row of zeros records that row~3 of $A$ was a combination of rows~1 and~2; elimination
knocked it out.
\end{example}

\begin{definition}[Echelon form, pivots, rank]\label{def:u08-rank}
A matrix is in \emph{echelon form} if every zero row lies below all nonzero rows and the first
nonzero entry of each nonzero row (its \emph{pivot}) lies strictly to the right of the pivot in
the row above. Elimination brings every matrix $A$ to an echelon form $U$. The \emph{rank} of
$A$, written $r=\rank A$, is the number of pivots. Columns that contain a pivot are
\emph{pivot columns}; the remaining columns are \emph{free columns}.
\end{definition}

A row holds at most one pivot and so does a column, so always
\[
r\le m\qquad\text{and}\qquad r\le n.
\]
In Example~\ref{ex:u08-elim} the rank is $2$, the pivot columns are $1$ and $3$, and the free
columns are $2$ and $4$. (That the rank does not depend on the order of the elimination steps
follows from Section~\ref{sec:u10}, where $r$ is identified with the dimension of the column
space, a number defined without elimination.)

\subsection{Pivot variables, free variables and special solutions}

The system $U\vb{x}=\vb{0}$ of Example~\ref{ex:u08-elim} is
\begin{align*}
x_1+2x_2+2x_3+2x_4&=0,\\
2x_3+4x_4&=0.
\end{align*}
The unknowns that multiply pivot columns, $x_1$ and $x_3$, are \emph{pivot variables}; the
others, $x_2$ and $x_4$, are \emph{free variables}. The free variables may be given any values.
Once they are chosen, back substitution determines the pivot variables uniquely, because each
nonzero row, read from the bottom up, contains exactly one new pivot variable with a nonzero
coefficient.

\begin{definition}[Special solutions]\label{def:u08-special}
For each free variable, the \emph{special solution} is the solution of $A\vb{x}=\vb{0}$ in
which that free variable equals $1$, every other free variable equals $0$, and the pivot
variables are found by back substitution.
\end{definition}

\begin{example}[The two special solutions]\label{ex:u08-special}
Take $x_2=1$, $x_4=0$. The second equation gives $2x_3=0$, so $x_3=0$; the first gives
$x_1+2=0$, so $x_1=-2$. Next take $x_2=0$, $x_4=1$. Then $2x_3+4=0$ gives $x_3=-2$, and
$x_1+2(-2)+2=0$ gives $x_1=2$:
\[
\vb{s}_1=\begin{bmatrix}-2\\1\\0\\0\end{bmatrix},\qquad
\vb{s}_2=\begin{bmatrix}2\\0\\-2\\1\end{bmatrix}.
\]
Each special solution is a dependency among the columns of $A$. The first says
$-2(\text{col }1)+(\text{col }2)=\vb{0}$: column~2 is twice column~1. The second says
$2(\text{col }1)-2(\text{col }3)+(\text{col }4)=\vb{0}$, and indeed
$2(1,2,3)-2(2,6,8)+(2,8,10)=(0,0,0)$. The whole nullspace is
\[
\Nul{A}=\{\,c\,\vb{s}_1+d\,\vb{s}_2 : c,d\in\R\,\},
\]
a two-dimensional plane through the origin in $\R^4$.
\end{example}

\begin{theorem}[Special solutions describe the nullspace]\label{thm:u08-special}
Let $A$ be $m\times n$ with rank $r$. There are exactly $n-r$ special solutions
$\vb{s}_1,\dots,\vb{s}_{n-r}$, and $\Nul{A}$ consists of all their linear combinations. More
precisely, if $\vb{x}\in\Nul{A}$ and its free variables have the values $c_1,\dots,c_{n-r}$,
then $\vb{x}=c_1\vb{s}_1+\dots+c_{n-r}\vb{s}_{n-r}$.
\end{theorem}
\begin{proof}
There are $r$ pivot columns, hence $n-r$ free columns and one special solution for each. The
special solutions lie in $\Nul{A}$, which is a subspace, so all their combinations do too.

Conversely, let $\vb{x}\in\Nul{A}$ have free variables $c_1,\dots,c_{n-r}$ and put
$\vb{y}=\vb{x}-\sum_j c_j\vb{s}_j$. Then $\vb{y}\in\Nul{A}=\Nul{U}$ by
Proposition~\ref{prop:u08-same-null}. Since $\vb{s}_j$ has a $1$ in the $j$th free position and
$0$ in the other free positions, every free variable of $\vb{y}$ is $0$. Now solve
$U\vb{y}=\vb{0}$ from the last nonzero row upward. In that row every entry to the right of the
pivot multiplies a free variable, which is zero, so the pivot variable is zero. In each row above,
the entries to the right of the pivot multiply free variables or pivot variables already shown to
be zero, so that row's pivot variable is zero as well. Hence $\vb{y}=\vb{0}$, that is,
$\vb{x}=\sum_j c_j\vb{s}_j$.
\end{proof}

\begin{corollary}[More unknowns than equations]\label{cor:u08-wide}
If $n>m$, then $A\vb{x}=\vb{0}$ has a solution $\vb{x}\ne\vb{0}$.
\end{corollary}
\begin{proof}
Since $r\le m<n$, there is at least one free variable, and its special solution has a $1$ in
that position.
\end{proof}

\subsection{The reduced row echelon form \texorpdfstring{\(R\)}{R}}

Elimination can go further. Working upward, we can also clear the entries \emph{above} each pivot
and then divide each pivot row by its pivot.

\begin{definition}[Reduced row echelon form]\label{def:u08-rref}
A matrix is in \emph{reduced row echelon form} if it is in echelon form, every pivot equals $1$,
and every pivot is the only nonzero entry in its column. Elimination brings every matrix $A$ to
such a matrix, and (a standard fact we use without proof) the result does not depend on the order
of the steps; it is denoted $R=\operatorname{rref}(A)$ (MATLAB: \texttt{rref(A)}).
\end{definition}

\begin{example}[From \texorpdfstring{\(U\)}{U} to \texorpdfstring{\(R\)}{R}]\label{ex:u08-rref}
In $U$ from Example~\ref{ex:u08-elim}, subtract row~2 from row~1 to clear the $2$ above the
second pivot, then divide row~2 by its pivot $2$:
\[
U=\begin{bmatrix}1&2&2&2\\0&0&2&4\\0&0&0&0\end{bmatrix}\longrightarrow
\begin{bmatrix}1&2&0&-2\\0&0&2&4\\0&0&0&0\end{bmatrix}\longrightarrow
R=\begin{bmatrix}1&2&0&-2\\0&0&1&2\\0&0&0&0\end{bmatrix}.
\]
The equations $R\vb{x}=\vb{0}$ are $x_1=-2x_2+2x_4$ and $x_3=-2x_4$: the pivot variables are
written directly in terms of the free ones. The free columns of $R$ contain $(2,0)$ and
$(-2,2)$ in the pivot rows, and these numbers reappear in $\vb{s}_1$ and $\vb{s}_2$ with their
signs reversed: $(x_1,x_3)=(-2,0)$ and $(2,-2)$.
\end{example}

The sign reversal is no accident. Suppose for the moment that the pivot columns come first (this
is only a renaming of the unknowns). Then $R$ has an identity block in its pivot rows and pivot
columns, a block $F$ in its pivot rows and free columns, and zero rows below:
\[
R=\begin{bmatrix}I&F\\0&0\end{bmatrix},\qquad I\text{ is }r\times r,\quad F\text{ is }r\times(n-r).
\]

\begin{proposition}[The nullspace matrix]\label{prop:u08-nullmatrix}
With $R$ in this block form, the $n\times(n-r)$ matrix
\[
N=\begin{bmatrix}-F\\ I\end{bmatrix}\qquad(\text{here }I\text{ is }(n-r)\times(n-r))
\]
satisfies $RN=0$, and its columns are the special solutions.
\end{proposition}
\begin{proof}
Block multiplication gives $RN=\begin{bmatrix}I(-F)+F\,I\\ 0\end{bmatrix}=0$. Column $j$ of $N$
has free part equal to column $j$ of the identity, a $1$ in the $j$th free position and zeros
in the others, so it is the solution whose free variables are those of the $j$th special
solution. Since back substitution determines the pivot variables uniquely, it \emph{is} the
$j$th special solution. Equivalently: $R\vb{x}=\vb{0}$ says
$\vb{x}_{\text{pivot}}+F\vb{x}_{\text{free}}=\vb{0}$, so $\vb{x}_{\text{pivot}}=-F\vb{x}_{\text{free}}$.
\end{proof}

When the pivot columns are not first, put the rows of $-F$ in the pivot positions and the rows
of $I$ in the free positions. In Example~\ref{ex:u08-rref}, $F=\begin{bmatrix}2&-2\\0&2\end{bmatrix}$
(rows~1--2, columns~2 and~4), so rows~1 and~3 of $N$ come from $-F$ and rows~2 and~4 from $I$:
\[
N=\begin{bmatrix}-2&2\\1&0\\0&-2\\0&1\end{bmatrix}=\begin{bmatrix}\vb{s}_1&\vb{s}_2\end{bmatrix}.
\]
Figure~\ref{fig:u08-blocks} records the shapes; this is exactly what the teaching code
\texttt{null(A)} does after computing $R$.

\begin{figure}[ht]\centering
\begin{tikzpicture}[scale=0.55,>=Stealth]
  % R: m = 4 rows (r = 2.5 + 1.5), n = 6 columns (r = 2.5 + 3.5)
  \draw[fill=blue!8] (0,1.5) rectangle (2.5,4);
  \draw[fill=orange!12] (2.5,1.5) rectangle (6,4);
  \draw[fill=gray!8] (0,0) rectangle (6,1.5);
  \node at (1.25,2.75) {$I$};
  \node at (4.25,2.75) {$F$};
  \node at (3,0.75) {$0$};
  \draw[decorate,decoration={brace,amplitude=4pt}] (0,4.15) -- (2.5,4.15)
    node[midway,above=4pt,font=\footnotesize,align=center] {$r$ pivot\\columns};
  \draw[decorate,decoration={brace,amplitude=4pt}] (2.5,4.15) -- (6,4.15)
    node[midway,above=4pt,font=\footnotesize,align=center] {$n-r$ free\\columns};
  \draw[decorate,decoration={brace,amplitude=4pt,mirror}] (-0.15,1.5) -- (-0.15,4)
    node[midway,left=5pt,font=\footnotesize,align=right] {$r$ pivot\\rows};
  \draw[decorate,decoration={brace,amplitude=4pt,mirror}] (-0.15,0) -- (-0.15,1.5)
    node[midway,left=5pt,font=\footnotesize,align=right] {$m-r$\\zero rows};
  \node at (3,-0.7) {$R$ \;($m\times n$)};
  % N: n = 6 rows (2.5 + 3.5), n - r = 3.5 columns
  \begin{scope}[xshift=8.6cm,yshift=-1cm]
    \draw[fill=orange!12] (0,3.5) rectangle (3.5,6);
    \draw[fill=blue!8] (0,0) rectangle (3.5,3.5);
    \node at (1.75,4.75) {$-F$};
    \node at (1.75,1.75) {$I$};
    \draw[decorate,decoration={brace,amplitude=4pt}] (3.65,6) -- (3.65,3.5)
      node[midway,right=5pt,font=\footnotesize,align=left] {pivot\\variables};
    \draw[decorate,decoration={brace,amplitude=4pt}] (3.65,3.5) -- (3.65,0)
      node[midway,right=5pt,font=\footnotesize,align=left] {free\\variables};
    \node at (1.75,-0.4) {$N$ \;($n\times(n-r)$)};
  \end{scope}
  \node at (7.3,2) {$\times$};
  \node at (15.6,2) {$=\;0$};
\end{tikzpicture}
\caption{Block shapes in $RN=0$ when the pivot columns come first. Each column of $N$ is a
special solution: the identity block sets one free variable to $1$ and the others to $0$, and
$-F$ supplies the pivot variables, so that $I(-F)+F\,I=0$.}
\label{fig:u08-blocks}
\end{figure}

\subsection{Two more examples}

\begin{example}[The transpose needs a row exchange]\label{ex:u08-transpose}
Take $A\T=\begin{bmatrix}1&2&3\\2&4&6\\2&6&8\\2&8&10\end{bmatrix}$. Its third column is the
sum of the first two, so we expect a free third column. Subtract $2\times$row~1 from rows 2, 3
and~4; column~2 then has a zero in the pivot position but a $2$ below it, so exchange rows~2
and~3; finally subtract $2\times$row~2 from row~4:
\[
\begin{bmatrix}1&2&3\\0&0&0\\0&2&2\\0&4&4\end{bmatrix}\longrightarrow
\begin{bmatrix}1&2&3\\0&2&2\\0&0&0\\0&4&4\end{bmatrix}\longrightarrow
U=\begin{bmatrix}1&2&3\\0&2&2\\0&0&0\\0&0&0\end{bmatrix}.
\]
Column~3 has nothing in the pivot position or below, so the rank is $2$, the same as the rank
of $A$; Section~\ref{sec:u11} proves that $\rank A\T=\rank A$ always. With the free variable
$x_3=1$, the second row gives $2x_2+2=0$, so $x_2=-1$, and the first gives $x_1-2+3=0$, so
$x_1=-1$. The nullspace is the line $\{c(-1,-1,1)\}$ in $\R^3$. Going on to
$R=\begin{bmatrix}1&0&1\\0&1&1\\0&0&0\\0&0&0\end{bmatrix}$ shows $F=\begin{bmatrix}1\\1\end{bmatrix}$
and $N=\begin{bmatrix}-F\\1\end{bmatrix}=(-1,-1,1)$, as it must.
\end{example}

\begin{example}[What a matrix crushes]\label{ex:u08-crush}
Let $A=\begin{bmatrix}1&2&2\\2&4&4\end{bmatrix}$, which takes $n=3$ inputs to $m=2$ outputs.
Row~2 minus $2\times$row~1 is $(0,0,0)$, so $R=\begin{bmatrix}1&2&2\\0&0&0\end{bmatrix}$ has one
pivot and $r=1$. The free columns $2$ and~$3$ give
\[
\vb{s}_1=\begin{bmatrix}-2\\1\\0\end{bmatrix},\qquad \vb{s}_2=\begin{bmatrix}-2\\0\\1\end{bmatrix},
\qquad A\vb{s}_1=\begin{bmatrix}-2+2\\-4+4\end{bmatrix}=\vb{0},\quad
A\vb{s}_2=\begin{bmatrix}-2+2\\-4+4\end{bmatrix}=\vb{0}.
\]
So $\Nul{A}$ is a plane in $\R^3$ of dimension $n-r=2$: a whole plane of inputs is sent to
zero. Both special solutions are perpendicular to the row $(1,2,2)$, since
$-2+2+0=0$ and $-2+0+2=0$. That is no coincidence: $A\vb{x}=\vb{0}$ says precisely that
$\vb{x}$ has zero dot product with every row. Section~\ref{sec:u11} completes this example with
the other three subspaces.
\end{example}

\begin{trap}
Row operations preserve the nullspace but \emph{not} the column space. In
Example~\ref{ex:u08-elim}, column~3 of $U$ is $(2,2,0)$. Every vector $A\vb{x}$ has third component
equal to the sum of the first two (row~3 of $A$ is row~1 plus row~2), and $2+2\ne 0$, so $(2,2,0)$
is not in $\Col{A}$. Use $U$ or $R$ to find $\Nul{A}$ and to locate the pivot columns, but take
those pivot columns from the original $A$ when you want the column space. A second slip: setting
\emph{all} free variables to zero produces only $\vb{x}=\vb{0}$; a special solution needs a $1$.
\end{trap}

\subsection{Exercises}

\begin{exercise}\label{exr:u08-rref}
Let $A=\begin{bmatrix}1&5&7&9\\0&4&1&7\\2&-2&11&-3\end{bmatrix}$. (a) Find the reduced row
echelon form of $A$. (b) What is the rank of $A$? (c) Find the special solutions of
$A\vb{x}=\vb{0}$.
\end{exercise}
\begin{solution}
(a) Subtract $2\times$row~1 from row~3, giving row~3 $=(0,-12,-3,-21)$. Divide row~2 by $4$:
$(0,1,\tfrac14,\tfrac74)$. Add $12\times$(new row~2) to row~3, giving a zero row. Subtract
$5\times$row~2 from row~1: $7-\tfrac54=\tfrac{23}{4}$ and $9-\tfrac{35}{4}=\tfrac14$. Hence
\[
R=\begin{bmatrix}1&0&\tfrac{23}{4}&\tfrac14\\[2pt]0&1&\tfrac14&\tfrac74\\[2pt]0&0&0&0\end{bmatrix}.
\]
(b) There are two pivots, so $\rank A=2$.
(c) The free variables are $x_3,x_4$, and $F=\begin{bmatrix}23/4&1/4\\1/4&7/4\end{bmatrix}$.
The columns of $\begin{bmatrix}-F\\I\end{bmatrix}$ are
$\vb{s}_1=\bigl(-\tfrac{23}{4},-\tfrac14,1,0\bigr)$ and $\vb{s}_2=\bigl(-\tfrac14,-\tfrac74,0,1\bigr)$.
Check one entry: row~2 of $A$ times $\vb{s}_1$ is $4(-\tfrac14)+1=0$.
\end{solution}

\begin{exercise}\label{exr:u08-change}
Change one entry of Example~\ref{ex:u08-crush}: $A=\begin{bmatrix}1&2&2\\2&4&5\end{bmatrix}$.
Find the rank, the special solutions and the dimension of the nullspace.
\end{exercise}
\begin{solution}
Row~2 minus $2\times$row~1 is now $(0,0,1)$, so there are pivots in columns~1 and~3 and $r=2$.
Clearing above the second pivot gives $R=\begin{bmatrix}1&2&0\\0&0&1\end{bmatrix}$. The only free
variable is $x_2$; setting $x_2=1$ gives $x_3=0$ and $x_1=-2$. The nullspace is the line through
$(-2,1,0)$, of dimension $3-2=1$: the plane of Example~\ref{ex:u08-crush} has shrunk to a line.
\end{solution}

\begin{exercise}\label{exr:u08-reverse}
Find a $2\times 4$ matrix $R$ in reduced row echelon form whose special solutions are
$(-3,1,0,0)$ and $(-2,0,-6,1)$.
\end{exercise}
\begin{solution}
In a special solution the free variables are $1$ and $0$; here the pattern
$(x_2,x_4)=(1,0)$ and $(0,1)$ shows that $x_2,x_4$ are free and $x_1,x_3$ are pivot variables.
By Proposition~\ref{prop:u08-nullmatrix} the pivot entries of the special solutions are the
columns of $-F$: $(x_1,x_3)=(-3,0)$ and $(-2,-6)$, so $F$ has columns $(3,0)$ and $(2,6)$.
Placing the identity in columns~1 and~3 and $F$ in columns~2 and~4,
\[
R=\begin{bmatrix}1&3&0&2\\0&0&1&6\end{bmatrix}.
\]
Check: $R(-3,1,0,0)=(-3+3,\,0)=\vb{0}$ and $R(-2,0,-6,1)=(-2+2,\,-6+6)=\vb{0}$.
\end{solution}

\begin{exercise}\label{exr:u08-rankAB}
Let $B=\begin{bmatrix}1&1\\1&1\end{bmatrix}$. Find $2\times2$ matrices $A_1$ and $A_2$ with
$\rank(A_1B)=1$ and $\rank(A_2B)=0$, and describe all possible $A_2$.
\end{exercise}
\begin{solution}
$A_1=I$ works: $A_1B=B$ has equal rows, so elimination leaves one pivot. $A_2=0$ works
trivially. In general, a row $(a,b)$ of $A_2$ multiplies $B$ to give $(a+b,\,a+b)$, which is zero
exactly when $b=-a$. So $A_2B=0$ precisely when every row of $A_2$ is a multiple of $(1,-1)$, for
example $A_2=\begin{bmatrix}1&-1\\1&-1\end{bmatrix}$.
\end{solution}
